\documentclass[conference]{IEEEtran}

\usepackage{amsmath,amssymb,bm}
\usepackage{booktabs}
\usepackage{cite}
\usepackage{url}

\newcommand{\C}{\mathbb{C}}
\newcommand{\Ftwo}{\mathbb{F}_2}
\newcommand{\norm}[1]{\left\lVert #1\right\rVert}
\newcommand{\Pset}{\mathcal{P}}
\newcommand{\Dset}{\mathcal{D}}
\newcommand{\Kset}{\mathcal{K}}

\title{JUNA-Lite: Posterior-Anchor Partial-FFT Combining and LDPC Decoding}

\author{\IEEEauthorblockN{Gabriel Chua Yu Han}}

\begin{document}
\maketitle

\begin{abstract}
Underwater acoustic communication channels are time varying and frequency
selective.  Residual Doppler destroys subcarrier orthogonality and produces
intercarrier interference.  Partial-FFT reception keeps short-time
observations, and a complex combining vector forms one pre-decoder soft symbol
per carrier.  A separate receiver estimates the combining vectors from pilots
and performs low-density parity-check decoding afterward.  JUNA-Lite closes
this information loop.  It starts from the same pilot-trained Partial-FFT
candidate, maps belief-propagation posterior metrics to soft QPSK symbols and
confidence values, combines reliable data symbols with the pilots, updates only the
combining vectors, and decodes again.  JUNA-Lite does not jointly optimize
the physical response model and the bit latents of the full JUNA objective.
The method is therefore a small posterior-anchor refinement of ratio
combining, not the full JUNA solver.  Joint use of pilot observations
and decoder information can improve the pre-decoder statistic relative to
separate combining and decoding, while confidence selection, pilot anchoring,
candidate comparison, and early stopping limit unreliable feedback.
\end{abstract}

\begin{IEEEkeywords}
underwater acoustic communication, OFDM, Doppler, partial FFT, LDPC, belief
propagation, ratio combining, iterative receiver
\end{IEEEkeywords}

\section{Introduction}

Orthogonal frequency-division multiplexing (OFDM) permits frequency-domain
equalization, but residual Doppler makes the channel vary within one OFDM
interval.  The frequency-domain operator is then no longer diagonal:
carrier-to-carrier leakage produces intercarrier interference (ICI), and a
single full-interval FFT collapses time-local information.  Partial-FFT
demodulation divides the useful interval into several segments and keeps one
observation from every segment~\cite{yerramalli}.

The RLS-based Partial-FFT Combiner (RPC) estimates a complex combining vector
from known pilots and forms one scalar soft symbol for the channel
decoder.  This is a strong separate receiver, but its information flow is one
way: pilots determine the front end, and low-density parity-check (LDPC)
decoding follows after soft-symbol formation.  The parity constraints do not
participate in the combiner fit.

Turbo equalization established the value of exchanging information between
adaptive equalization and channel decoding~\cite{laot,zheng}.  JUNA-Lite
applies the same principle to Partial-FFT combining.  The decoder posterior
means become additional soft data anchors, and the decoder confidence
determines their weight.  The resulting receiver repeats ratio combining
and decoding without introducing the full branch-domain model of the general
JUNA formulation.

We present a compact mathematical description of this receiver using
the notation of the JUNA system model.  The central comparison is precise:
the separate receiver solves a pilot-only ridge problem and decodes once;
JUNA-Lite starts from that solution, adds confidence-weighted posterior
anchors, updates \(\mathbf W\), re-decodes, and keeps the better candidate.

\section{OFDM Under Residual Doppler}

Let \(N\) be the number of OFDM samples after cyclic-prefix removal, let
\(\mathbf r\in\C^N\) be the received block, and let
\(\mathbf F\in\C^{N\times N}\) be the unitary FFT matrix.  The active-carrier
set is \(\Kset\), the pilot set is \(\Pset\), and the data set is
\(\Dset=\Kset\setminus\Pset\).  The transmitted symbol on carrier \(k\) is
\(X_k\), the vector \(\mathbf X\in\C^N\) collects these symbols, and
\(\mathbf F^H\mathbf X\) is the transmitted block.

With an adequate cyclic prefix and a channel that does not vary over the
block, cyclic-prefix removal leaves a circulant operator.  The FFT
diagonalizes it, and every carrier obeys the one-tap model
\begin{equation}
  Y_k^{\mathrm{full}}
  =
  H_{k,k}X_k+\nu_k,
  \label{eq:one-tap}
\end{equation}
where \(Y_k^{\mathrm{full}}\) is the full-interval FFT bin, \(H_{k,k}\) is the
response on carrier \(k\), and \(\nu_k\) is noise.  Pilot-linear equalization
estimates the response on the pilots and interpolates it over \(\Kset\).

Motion time-scales the received packet.  After coarse correction a residual
scale \(a\) and a common frequency offset \(\theta\) remain, so carrier
\(\ell\) arrives at index \((1+a)\ell+\theta\) rather than at \(\ell\).  Bin
\(k\) multiplies by \(e^{-j2\pi kn/N}\) and averages, so carrier \(\ell\)
enters bin \(k\) with the coefficient
\begin{align}
  H_{k,\ell}
  &=
  \frac{1}{N}\sum_{n=0}^{N-1}
  e^{j2\pi\Delta_{k,\ell}n/N},
  \label{eq:ici-coefficient}\\
  \Delta_{k,\ell}
  &=
  (\ell-k)+\theta+a\ell,
  \label{eq:leftover-turns}
\end{align}
where \(\Delta_{k,\ell}\) is the number of turns carrier \(\ell\) makes
relative to bin \(k\) over one block.  When \(a\) and \(\theta\) are zero, the
turn count is an integer, the sum closes, and the coefficient vanishes for
every \(\ell\neq k\).  This is subcarrier orthogonality.  A nonzero \(a\) or
\(\theta\) makes the turn count fractional, the sum no longer closes, and the
residue is the leakage.

We consider \(N=1024\), \(k=4\), \(\ell=5\), \(a=0.007\), and \(\theta=0.20\).
Carrier \(5\) arrives at index \(5.235\), bin \(4\) sees
\(\Delta_{4,5}=1.235\) turns, and \eqref{eq:ici-coefficient} gives
\(H_{4,5}\approx0.129+0.116j\) with \(|H_{4,5}|\approx0.173\).  Without
Doppler the same pair gives \(\Delta_{4,5}=1\) and \(H_{4,5}=0\).

Collecting the coefficients of all active carriers,
\begin{equation}
  Y_k^{\mathrm{full}}
  =
  H_{k,k}X_k
  +
  \sum_{\ell\in\Kset\setminus\{k\}}
  H_{k,\ell}X_\ell
  +
  \nu_k,
  \label{eq:ici-row}
\end{equation}
where the first term carries the wanted symbol and the sum is the ICI.  The
magnitude of \eqref{eq:ici-coefficient} is a Dirichlet kernel in
\(\Delta_{k,\ell}\), so for small \(a\) and \(\theta\) the leakage stays near
the diagonal and the operator is banded.  Read in the frequency domain, the
same coefficient is the spectrum of carrier \(\ell\) sampled away from its
peak.  The bin grid no longer lands on the zeros of the neighboring carriers,
and it samples their tails instead.

Banded-ICI and sparse-channel receivers estimate the coefficients
\(H_{k,\ell}\), or the path parameters that produce them.  Partial-FFT
reception does not.  It keeps time-local observations of the same
full-interval statistic and forms one soft symbol from them by combining.
The response \(\mathbf C\) of the full JUNA objective is the Partial-FFT
counterpart of the banded matrix that \(H_{k,\ell}\) forms; JUNA-Lite
estimates neither, and updates \(\mathbf W\) alone.

\section{Partial-FFT Observation Model}

The full-interval statistic of \eqref{eq:ici-row} is one number per carrier.
Partial-FFT demodulation keeps several, one from each segment of the useful
interval.
Let \(g_m[n]\), \(m=1,\ldots,M\), be deterministic non-overlapping segment
windows, and define
\begin{equation}
  \mathbf S_m
  =
  \operatorname{diag}\!\left(g_m[0],\ldots,g_m[N-1]\right).
  \label{eq:segment}
\end{equation}
The branch observation on carrier \(k\), branch \(m\), is
\begin{align}
  Y_{k,m}
  &=
  \mathbf e_k^{T}\mathbf F\mathbf S_m\mathbf r,
  \label{eq:branch-operator}\\
  &=
  \frac{1}{\sqrt N}\sum_{n=0}^{N-1}
  g_m[n]r[n]e^{-j2\pi kn/N}.
  \label{eq:branch-sum}
\end{align}
The Partial-FFT observation vector is
\begin{equation}
  \mathbf Y_k
  =
  [Y_{k,1},\ldots,Y_{k,M}]^T\in\C^M .
  \label{eq:observation-vector}
\end{equation}
When the windows form a partition of unity,
\(\sum_{m=1}^{M}g_m[n]=1\), the branch sum equals the full FFT:
\begin{equation}
  \sum_{m=1}^{M}Y_{k,m}
  =
  \mathbf e_k^T\mathbf F
  \left(\sum_{m=1}^{M}\mathbf S_m\right)\mathbf r
  =
  Y_k^{\mathrm{full}}.
  \label{eq:branch-identity}
\end{equation}
Thus, Partial-FFT reception does not create another measurement; it keeps
temporally localized components of the full statistic.

\section{Separate Ratio Combining and Decoding}

Using the notation of the JUNA model, the pre-decoder scalar soft symbol is
\begin{equation}
  \widetilde X_k
  =
  \mathbf W_k^H\mathbf Y_k,
  \qquad
  \mathbf W_k\in\C^M .
  \label{eq:scalar-combiner}
\end{equation}
For a target carrier \(k\), let \(\Pset(k)\subseteq\Pset\) be a nearby pilot
set.  The separate pilot-only ridge criterion is
\begin{equation}
  J_k(\mathbf W_k)
  =
  \sum_{p\in\Pset(k)}
  \left|\mathbf W_k^H\mathbf Y_p-X_p\right|^2
  +
  \delta\norm{\mathbf W_k}_2^2,
  \qquad \delta>0 .
  \label{eq:rpc-objective}
\end{equation}
With
\begin{align}
  \mathbf R_k
  &=
  \sum_{p\in\Pset(k)}
  \mathbf Y_p\mathbf Y_p^H+\delta\mathbf I,
  \label{eq:rpc-correlation}\\
  \mathbf r_k
  &=
  \sum_{p\in\Pset(k)}
  \mathbf Y_pX_p^*,
  \label{eq:rpc-cross}
\end{align}
the normal equations and solution are
\begin{align}
  \mathbf R_k\mathbf W_k
  &=
  \mathbf r_k,
  \label{eq:rpc-normal}\\
  \widehat{\mathbf W}_k
  &=
  \mathbf R_k^{-1}\mathbf r_k.
  \label{eq:rpc-solution}
\end{align}

The implementation uses one vector for each carrier band.  If
\(\Kset_b\subseteq\Kset\) is band \(b\), then
\begin{equation}
  \mathbf W_k=\mathbf W_b,
  \qquad k\in\Kset_b .
  \label{eq:band-sharing}
\end{equation}
When a band has fewer than \(M\) local anchors, all available anchors are used.
After combining, a residual pilot response is
\begin{equation}
  \widehat H_p
  =
  \frac{\widetilde X_p}{X_p},
  \qquad p\in\Pset .
  \label{eq:pilot-response}
\end{equation}
Interpolation gives \(\widehat H_k\) on active carriers, and the equalized
symbol is
\begin{equation}
  \widehat X_k
  =
  \frac{\widetilde X_k}{\widehat H_k}.
  \label{eq:residual-equalization}
\end{equation}

For QPSK, let \(\widehat{\sigma}^2\) be the pilot mean-squared error with a
positive floor.  The signed channel metrics for the in-phase and quadrature
bits are
\begin{align}
  L_{i_1(k)}
  &=
  -\frac{2\operatorname{Re}\{\widehat X_k\}}
  {\widehat{\sigma}^2},
  \label{eq:qpsk-metric-real}\\
  L_{i_2(k)}
  &=
  -\frac{2\operatorname{Im}\{\widehat X_k\}}
  {\widehat{\sigma}^2}.
  \label{eq:qpsk-metric-imag}
\end{align}
The metrics are limited before belief-propagation decoding.  The separate
receiver performs this chain once:
\[
  \mathbf Y_k
  \longrightarrow
  \widehat{\mathbf W}_k
  \longrightarrow
  \widehat X_k
  \longrightarrow
  \{L_i\}
  \longrightarrow
  \widehat{\mathbf b}.
\]

\section{LDPC Posterior Symbols}

An LDPC code is specified by a sparse parity-check matrix
\(\mathbf P\in\Ftwo^{m_c\times n_b}\).  A codeword
\(\mathbf b\in\Ftwo^{n_b}\) satisfies
\begin{equation}
  \mathbf P\mathbf b^T=\mathbf 0
  \qquad\text{over }\Ftwo .
  \label{eq:parity}
\end{equation}
Let \(\mathcal N_j\) be the variable-node neighborhood of check \(j\), and let
\(\mathcal M_i\) be the check-node neighborhood of coded bit \(i\).  With
channel log-likelihood ratio \(\Lambda_i^{\mathrm{ch}}\), the variable-to-check
message is
\begin{equation}
  q_{i\rightarrow j}^{(\ell)}
  =
  \Lambda_i^{\mathrm{ch}}
  +
  \sum_{j'\in\mathcal M_i\setminus j}
  r_{j'\rightarrow i}^{(\ell-1)}.
  \label{eq:variable-message}
\end{equation}
The measured receiver uses normalized min-sum check updates:
\begin{equation}
  r_{j\rightarrow i}^{(\ell)}
  =
  \alpha
  \left[
  \prod_{i'\in\mathcal N_j\setminus i}
  \operatorname{sign}
  \left(q_{i'\rightarrow j}^{(\ell)}\right)
  \right]
  \min_{i'\in\mathcal N_j\setminus i}
  \left|q_{i'\rightarrow j}^{(\ell)}\right|,
  \label{eq:check-message}
\end{equation}
where \(\alpha=0.8\).  The posterior log-likelihood ratio and hard bit are
\begin{align}
  \Lambda_i^{\mathrm{post},(\ell)}
  &=
  \Lambda_i^{\mathrm{ch}}
  +
  \sum_{j\in\mathcal M_i}
  r_{j\rightarrow i}^{(\ell)},
  \label{eq:posterior-llr}\\
  \widehat b_i^{(\ell)}
  &=
  \begin{cases}
    0, & \Lambda_i^{\mathrm{post},(\ell)}\geq0,\\
    1, & \Lambda_i^{\mathrm{post},(\ell)}<0.
  \end{cases}
  \label{eq:hard-bit}
\end{align}
Decoding stops when \(\mathbf P\widehat{\mathbf b}^{T}=\mathbf0\) or the BP
iteration limit is reached.  The implementation returns the opposite metric
sign,
\begin{equation}
  L_i^{\mathrm{post},(t)}
  =
  -\Lambda_i^{\mathrm{post},(t)}.
  \label{eq:returned-posterior}
\end{equation}
Using the returned BP posterior metric for coded bit \(i\) at iteration \(t\),
the bit-latent notation of the JUNA model is
\begin{align}
  z_i^{(t)}
  &=
  -\frac{1}{2}L_i^{\mathrm{post},(t)},
  \label{eq:posterior-latent}\\
  \xi_i(\mathbf z^{(t)})
  &=
  \tanh z_i^{(t)}
  =
  -\tanh\!\left(\frac{L_i^{\mathrm{post},(t)}}{2}\right).
  \label{eq:posterior-bit}
\end{align}
For Gray-labeled QPSK, data carrier \(k\) carries coded bits \(i_1(k)\) and
\(i_2(k)\), and its posterior-mean symbol is
\begin{equation}
  s_k(\mathbf z^{(t)})
  =
  \frac{1}{\sqrt 2}
  \left[
  \xi_{i_1(k)}(\mathbf z^{(t)})
  +
  j\xi_{i_2(k)}(\mathbf z^{(t)})
  \right].
  \label{eq:posterior-symbol}
\end{equation}
The data-anchor confidence is the smaller bipolar magnitude:
\begin{equation}
  \omega_k^{(t)}
  =
  \min\!\left(
  \left|\xi_{i_1(k)}(\mathbf z^{(t)})\right|,
  \left|\xi_{i_2(k)}(\mathbf z^{(t)})\right|
  \right),
  \qquad 0\leq\omega_k^{(t)}\leq1.
  \label{eq:confidence}
\end{equation}
For BPSK, the corresponding definitions are
\begin{align}
  s_k(\mathbf z^{(t)})
  &=
  \xi_{i(k)}(\mathbf z^{(t)}),
  \label{eq:bpsk-symbol}\\
  \omega_k^{(t)}
  &=
  \left|\xi_{i(k)}(\mathbf z^{(t)})\right|.
  \label{eq:bpsk-confidence}
\end{align}

\section{JUNA-Lite Posterior-Anchor Estimation}

Let \(\tau\geq0\) be the confidence threshold and let \(K_{\max}\) be the
maximum number of data anchors.  The selected set is
\begin{equation}
  \mathcal Q^{(t)}
  =
  \left\{
  k\in\Dset:
  \omega_k^{(t)}\geq\tau
  \right\}.
  \label{eq:eligible}
\end{equation}
If \(|\mathcal Q^{(t)}|>K_{\max}\), only the \(K_{\max}\) largest confidence
values are retained.  In the deployed receiver \(\tau=0\) and \(K_{\max}\)
does not bind, so every decoded data carrier becomes an anchor and the
confidence weights alone limit unreliable feedback.  The complete target set
is
\begin{equation}
  \mathcal T^{(t)}
  =
  \Pset\cup\mathcal Q^{(t)}.
  \label{eq:target-set}
\end{equation}
We define the target and target weight by
\begin{align}
  d_k^{(t)}
  &=
  \begin{cases}
    X_k, & k\in\Pset,\\
    s_k(\mathbf z^{(t)}), & k\in\mathcal Q^{(t)},
  \end{cases}
  \label{eq:target}\\
  v_k^{(t)}
  &=
  \begin{cases}
    1, & k\in\Pset,\\
    \omega_k^{(t)}, & k\in\mathcal Q^{(t)}.
  \end{cases}
  \label{eq:target-weight}
\end{align}

For band \(\Kset_b\), the JUNA-Lite combiner update is the
confidence-weighted ridge problem
\begin{equation}
  \begin{aligned}
    \widehat{\mathbf W}_b^{(t+1)}
    =
    \arg\min_{\mathbf W_b\in\C^M}
    \Bigg\{&
    \sum_{k\in\mathcal T_b^{(t)}}
    v_k^{(t)}
    \left|
    \mathbf W_b^H\mathbf Y_k-d_k^{(t)}
    \right|^2\\
    &+
    \delta\norm{\mathbf W_b}_2^2
    \Bigg\}.
  \end{aligned}
  \label{eq:lite-objective}
\end{equation}
where
\begin{equation}
  \mathcal T_b^{(t)}
  =
  \mathcal T^{(t)}\cap\Kset_b .
  \label{eq:band-targets}
\end{equation}
The weighted correlation matrix and cross-correlation vector are
\begin{align}
  \mathbf R_b^{(t)}
  &=
  \sum_{k\in\mathcal T_b^{(t)}}
  v_k^{(t)}
  \mathbf Y_k\mathbf Y_k^H
  +
  \delta\mathbf I,
  \label{eq:lite-correlation}\\
  \mathbf r_b^{(t)}
  &=
  \sum_{k\in\mathcal T_b^{(t)}}
  v_k^{(t)}
  \mathbf Y_k
  \left(d_k^{(t)}\right)^* .
  \label{eq:lite-cross}
\end{align}
Therefore
\begin{align}
  \mathbf R_b^{(t)}
  \widehat{\mathbf W}_b^{(t+1)}
  &=
  \mathbf r_b^{(t)},
  \label{eq:lite-normal}\\
  \widehat{\mathbf W}_b^{(t+1)}
  &=
  \left(\mathbf R_b^{(t)}\right)^{-1}
  \mathbf r_b^{(t)}.
  \label{eq:lite-solution}
\end{align}
The new scalar symbol for \(k\in\Kset_b\) is
\begin{equation}
  \widetilde X_k^{(t+1)}
  =
  \left(\widehat{\mathbf W}_b^{(t+1)}\right)^H
  \mathbf Y_k.
  \label{eq:lite-output}
\end{equation}
Residual pilot equalization, metric formation, and BP decoding then produce
the next posterior metrics \(L_i^{\mathrm{post},(t+1)}\).

Equations~\eqref{eq:lite-objective}--\eqref{eq:lite-solution} are the complete
JUNA-Lite optimization step.  The update estimates \(\mathbf W\) only.  It does
not form or optimize the branch-domain response \(\mathbf C\), and it does not
minimize the full JUNA objective over
\((\mathbf C,\mathbf W,\mathbf z)\).  The variables \(\mathbf z^{(t)}\) in
JUNA-Lite are a representation of BP posterior metrics, not free variables in
a joint gradient solve.

\section{Candidate Selection and Stopping}

Let candidate \(c\) have parity validity \(q_c\in\{0,1\}\), syndrome weight
\(\sigma_c\), pilot error \(E_{\mathrm{pil},c}\), posterior tie error \(E_c\),
decoder score \(D_c\), and mean posterior magnitude \(\mu_c\).  These
quantities are
\begin{align}
  E_{\mathrm{pil},c}
  &=
  \frac{1}{|\Pset|}
  \sum_{p\in\Pset}
  \left|\widehat X_p-X_p\right|^2,
  \label{eq:pilot-error}\\
  \mu_c
  &=
  \frac{1}{n_b}
  \sum_{i=1}^{n_b}
  \left|L_i^{\mathrm{post}}\right|,
  \label{eq:posterior-magnitude}
\end{align}
and the posterior tie error is
\begin{equation}
  E_c
  =
  \frac{
  \displaystyle
  \sum_{k\in\Dset}
  \max(\omega_k,\epsilon)
  \left|
  \widehat X_k-s_k(\mathbf z)
  \right|^2
  }{
  \displaystyle
  \sum_{k\in\Dset}
  \max(\omega_k,\epsilon)
  }.
  \label{eq:tie-error}
\end{equation}
The decoder score is
\begin{equation}
  D_c
  =
  E_{\mathrm{pil},c}
  +0.25E_c
  +0.05\frac{\sigma_c}{m_c}
  -10^{-4}\mu_c.
  \label{eq:decoder-score}
\end{equation}
In the deployed receiver \(E_{\mathrm{pil},c}\) is inactive and takes the same
value for every candidate.  The residual pilot equalization of
\eqref{eq:pilot-response}--\eqref{eq:residual-equalization} divides each pilot
carrier by its measured response, so \(\widehat X_p=X_p\) to numerical
precision.  Hence the score separates candidates through \(E_c\), the
normalized syndrome weight, and \(\mu_c\).  We retain
\(E_{\mathrm{pil},c}\) because it separates candidates when a receiver scores
the combiner output \(\widetilde X_p\) before residual pilot equalization.
Candidate comparison is ordered in the implementation:
\begin{equation}
  \begin{aligned}
    c_1\succ c_0
    \quad\text{if the first true line is}\quad&\\
    q_{c_1}>q_{c_0},\qquad&\\
    \sigma_{c_1}<\sigma_{c_0},\qquad&\\
    D_{c_1}<D_{c_0}
    -0.005\max(|D_{c_0}|,\epsilon),\qquad&\\
    \mu_{c_1}>1.01\mu_{c_0}+10^{-6}.&
  \end{aligned}
  \label{eq:candidate-order}
\end{equation}
The first satisfied strict comparison determines the better candidate.  The
receiver starts with the pilot-trained Partial-FFT candidate.  A valid seed is
returned immediately.  Otherwise, posterior-anchor estimation and re-decoding are
repeated until a valid parity check, no improvement, or the iteration limit.
The best candidate observed during the iterations is returned.

\section{Joint and Separate Processing}

The separate receiver uses only \(\Pset\) in
\eqref{eq:rpc-objective}.  JUNA-Lite uses
\(\Pset\cup\mathcal Q^{(t)}\) in \eqref{eq:lite-objective}.  The two methods
use the same Partial-FFT observations \(\mathbf Y_k\), the same pilots, the
same ridge term, the same residual pilot equalization, and the same BP decoder.
Their difference is the information flow.

The separate chain is
\[
  \{\mathbf Y_k\}
  \rightarrow
  \mathbf W
  \rightarrow
  \{\widehat X_k\}
  \rightarrow
  \{L_i\}
  \rightarrow
  \widehat{\mathbf b}.
\]
The JUNA-Lite chain is
\[
  \{\mathbf Y_k\}
  \rightarrow
  \mathbf W^{(t)}
  \rightarrow
  \{L_i^{\mathrm{post},(t)}\}
  \rightarrow
  \{s_k(\mathbf z^{(t)}),\omega_k^{(t)}\}
  \rightarrow
  \mathbf W^{(t+1)}.
\]
Better combining can produce better decoder information; better decoder
information can provide better anchors; and better anchors can produce better
combining vectors.  This feedback is useful when pilots alone do not specify
a stable combiner and when the decoder posterior information is reliable.  A
valid initial candidate produces no refinement, and unreliable feedback can
reduce the value of the additional anchors.  Confidence selection, fixed
pilot anchors, ridge regularization, candidate comparison, and early stopping
are therefore essential parts of the method.

\section{Complexity and Evaluation}

For each band and iteration, formation of
\(\mathbf R_b^{(t)}\) requires
\(\mathcal O(|\mathcal T_b^{(t)}|M^2)\) operations, and the small linear solve
requires \(\mathcal O(M^3)\).  Application of the combiner to all active
carriers requires \(\mathcal O(|\Kset|M)\).  BP decoding is performed once for
the seed and once for every accepted refinement step.  JUNA-Lite avoids the
branch-domain \(\mathbf C\) update and the nonlinear \(\mathbf z\) optimization
of the full JUNA solver.

We require identical modulation, code, pilot pattern,
Partial-FFT branches, channel, signal-to-noise ratio, and packet budget.  We
compare pilot-only Partial-FFT plus BP against JUNA-Lite
initialized from the same candidate.  Relevant measures include coded bit
error rate, frame error rate, valid parity checks, pilot overhead, processing
latency, iteration count, and computational complexity.  A fixed-feedback
control uses pilots only inside the same iterative path; oracle and corrupted
feedback controls separate the value of decoder information from the other
receiver operations.

\section{Conclusion}

JUNA-Lite is a confidence-weighted posterior-anchor refinement of Partial-FFT
ratio combining.  It starts from the same pilot-trained combiner used by the
separate receiver, obtains BP posterior information, adds reliable data
symbols to the pilot anchors, estimates \(\mathbf W\), and decodes again.  The
method does not optimize \(\mathbf C\), and its \(\mathbf z\) variables are
posterior representations rather than free joint-solver variables.  Joint use
of the observations and decoder information can improve the pre-decoder
statistic when pilots alone do not specify a stable combiner.  The seed,
confidence, ridge, comparison, and stopping rules retain the separate
receiver as a fallback.

\begin{thebibliography}{00}

\bibitem{yerramalli}
S. Yerramalli, M. Stojanovic, and U. Mitra,
``Partial FFT Demodulation: A Detection Method for Highly Doppler Distorted
OFDM Systems,''
\emph{IEEE Transactions on Signal Processing}, 2012.

\bibitem{laot}
C. Laot, A. Glavieux, and J. Labat,
``Turbo Equalization: Adaptive Equalization and Channel Decoding Jointly
Optimized,''
\emph{IEEE Transactions on Communications}, 2001.

\bibitem{zheng}
Y. R. Zheng, J. Wu, and C. Xiao,
``Turbo Equalization for Single-Carrier Underwater Acoustic Communications,''
\emph{IEEE Communications Magazine}.

\bibitem{juna}
G. C. Y. Han,
\emph{JUNA: A Joint Receiver for Doppler-Resilient OFDM: Partial-FFT
Coefficient Estimation and LDPC-Coupled Detection}.

\end{thebibliography}

\end{document}
